\section{Who Decides, and the Power to Hurt}\label{sec:decide}

So far Europe has been one actor. In fact The Hague licenses ASML's shipments to China, ASML decides whether to make them, and the Union decides whether its tools can be turned against the United States (\cref{sec:background}). Europe can also impose costs on Washington while the two bargain (\cref{sec:holdup}).

\subsection{The gatekeeper}\label{sec:gatekeeper}

Exit toward China is a Dutch decision. The licence is granted where the exporter is established, and neither other member states nor the Commission can refuse it \citep{eu2021dualuse}. Let $Q_{NL}=\omega_{NL}-c_{NL}$ be the gatekeeper's own exit premium, net of the credible costs of exit that fall on the Netherlands. By \cref{lem:bidding}, Washington keeps the Netherlands aligned at price $\pos{Q_{NL}}$.

Washington can lower that price by divide and rule. A benefit that it provides to a loyal Netherlands in the baseline, and credibly withdraws if The Hague licenses exit, lowers $Q_{NL}$ one for one. It cuts Washington's payment by up to the benefit's value, and once the gatekeeper has slack, every further unit adds to what Washington can spend on other demands (\cref{rem:gatekeeper}, \cref{prop:slack}). A side payment promised on top of the baseline saves Washington nothing. On our reading the results of \cref{sec:slack,sec:control} carry over to China-bound licensing with $Q_{NL}$ in place of $Q$ and each instrument measured by what the Netherlands bears, but not to US-facing measures, which the Union decides.

European partners move $\omega_{NL}$ or $c_{NL}$ only through transfers: compensation, coordinated retaliation or cover against Chinese counter-coercion. Washington can tier member states directly. By our computation from 2026 Council population figures \citep{council2025population}, the ten members in the AI Diffusion Rule's favoured group (\cref{sec:background}) held 70.46 percent of EU population but too few members for a qualified majority, while the other seventeen held 29.54 percent and enough members to block. BIS withdrew the rule before compliance was due, so tiering is latent, but one rule can reward the gatekeeper and split the Council at once.

\subsection{The firm between two jurisdictions}\label{sec:firm}

A licence permits a shipment and orders nothing. ASML ships when its Chinese business, plus the penalties China ($p_C$) and the Union's Blocking Statute ($p_B$) would impose for complying with US measures instead, at least matches the credible US penalty for shipping. Otherwise the state must compel it at a cost $\zeta$, and Europe does so only if its premium exceeds that cost (\cref{rem:firm}).

Since April 2026, $p_C$ has a legal basis. Decree 835 lets Beijing blacklist, bar and expose to damages suits and fines those who implement foreign extraterritorial measures \citep{reedsmith2026decree}. ASML's staff in China who implemented a US servicing ban would be exposed. A higher $p_C$ makes a licensed firm more willing to ship and weakly raises Europe's effective premium, without touching the credibility of Washington's jurisdiction. For a US ban on servicing in China, the same penalty enters \cref{prop:counter}(a) as a counter-threat $r_k$ only if Europe internalizes the firm's loss.

The Blocking Statute's annex lists US sanctions on Cuba and Iran and no US export-control law \citep{eu1996blocking, ec2018blocking}, so $p_B=0$ today. Even a listing would be bounded: in \emph{Bank Melli} the Court of Justice let national courts weigh the firm's economic loss \citep{cjeu2021bankmelli}, which we read as keeping $p_B$ small for a firm with ASML's US business.

The Hague compelled Nexperia under the Goods Availability Act, China blocked exports from Nexperia's Chinese plants, and the order was suspended after Washington had suspended the rule behind its demand \citep{kst32637720, kst32637_717}. If the Chinese ban moved The Hague, $\zeta$ is high when China retaliates against compulsion itself, and the record does not tell us which (\cref{sec:record}).

\subsection{The Council}\label{sec:council}

Turning ASML against Washington is collective, because its most advanced tools ship to US fabs under the Union general export authorisation EU001 (\cref{sec:background}). Under the Anti-Coercion Instrument the Council first determines by qualified majority that coercion exists. Response measures are Commission implementing acts; if the examination committee gives no opinion, the act goes to an appeal committee, where the Commission may adopt unless a qualified majority votes against \citep{eu2023aci, eu2011comitology}. A delegated act removing the United States from EU001 would instead apply at once and stand unless the European Parliament or a reinforced qualified majority of the Council objected \citep{eu2021dualuse, eu2016tfeu}, but The Hague would still have to deny the individual licences.

With enforceable vote-contingent payments, Washington defeats a countermeasure most cheaply by buying a blocking minority before the Council's determination, since every later stage of either route needs a coalition at least as large (\cref{rem:council}). Because Washington can make posture conditional member by member (\cref{sec:security}), the members whose defence depends most on US forward presence join for free. They cannot block alone: the six frontline states hold 15.09 percent of the population, and a block by number needs thirteen eastern and Nordic members (\cref{tab:council}). We expect opposition to an anti-US measure to start on the frontline and to block only once it spreads to central Europe and the Nordics, or a large state joins. The seven states that \emph{Europe 2031} names as opposing or abstaining on its failed 2029 measure hold 28.06 percent of EU population \citep{juijn2026europe, council2025population} and cannot block. Any blocking completion needs Germany, France or Spain, or at least three further states (Belgium, Czechia and Portugal would do), or six more of any size to reach thirteen members (\cref{app:numerics}).

The formal threshold is not what binds. Recital 85 of the instrument commits the Commission, at every stage including appeal, to avoid ``going against any predominant position amongst Member States'' \citep{eu2023aci}, and its Statement C/2023/1341 repeats the commitment \citep{commission2023acistatement}. A countermeasure that large states oppose should die before any vote: no examination, no proposal, or a package without the ASML element. Member states gave positive committee opinions to rebalancing measures against US goods on 9 April and 24 July 2025 \citep{commission2025usstatement, commission2025rebalancing1564}. Neither applied. The April package was suspended after the US tariff pause and the July package after the Turnberry agreement \citep{commission2025suspension786, commission2025suspension1727}; the suspension was renewed in January 2026 after Washington dropped the tariffs it had threatened over Greenland \citep{dazio2026greenland, tbc2026summaryjan}, and again in July 2026 with no end date \citep{commission2026suspension1893}. The Union arms and then suspends on a US concession. No armed measure has been tested against a Washington that refused to move. Washington has a route cheaper than buying a blocking minority: a concession that triggers suspension.

\subsection{Hold-up}\label{sec:holdup}

\Cref{prop:master} lets the status quo run unchanged while Europe and Washington bargain. In a dispute over control of ASML or over Europe's access to frontier AI, Washington could cut Europe's AI allocation while talks continue, and Europe could interrupt ASML's servicing of US-aligned fabs. Let disagreement carry flow costs $c_E(t)$ to Europe and $c_U(t)$ to the United States, and write $\hat C_i(t)=(r+\lambda)\int_0^\infty e^{-(r+\lambda)s}c_i(t+s)\,ds$ for the average discounted cost from date $t$ on.

\begin{proposition}[Pain and rupture]\label{prop:pain}
Consider the game of \cref{prop:master} with $r>0$ and disagreement flow costs $c_E(t),c_U(t)\ge0$ that are bounded and piecewise uniformly continuous, with finitely many discontinuities in every bounded interval. Let $M=Q+\beta(V-Q)$, $n(t)=\beta c_U(t)-(1-\beta)c_E(t)$, $N(t)=\beta\hat C_U(t)-(1-\beta)\hat C_E(t)$ and
\[
x^\circ(t)=\psi M+(1-\psi)N(t).
\]
\begin{enumerate}[label=(\roman*)]
\item For every $\Delta>0$, Europe's Markov perfect equilibrium payoff is unique. In every round in which $\lambda>0$ or $c_E+c_U>0$, every Markov perfect equilibrium has immediate agreement. If $\lambda=0$ and $c_E(t)+c_U(t)=0$ on an interval, there are also payoff-equivalent equilibria with delay during that interval.
\item As $\Delta\to0$, Europe's equilibrium payoff converges to
\begin{align*}
P^*&=\sup_{\tau\in[0,\infty]}\Big\{\int_0^\tau e^{-(r+\lambda)s}\big(\lambda M+r\,n(s)\big)\,ds+e^{-(r+\lambda)\tau}Q\Big\}\\
&=x^\circ(0)+\pos{\sup_{t\ge0}e^{-(r+\lambda)t}\big(Q-x^\circ(t)\big)}.
\end{align*}
\item $P^*=x^\circ(0)=\psi\big[Q+\beta(V-Q)\big]+(1-\psi)\big[\beta\hat C_U(0)-(1-\beta)\hat C_E(0)\big]$ if and only if
\[
Q\le\inf_{t\ge0}x^\circ(t)\quad\Longleftrightarrow\quad Q\le\bar Q\equiv\inf_{t\ge0}\frac{\lambda\beta V+r\,N(t)}{r+\lambda\beta}.
\]
Under this condition the limiting equilibrium value at every date $t$ is $x^\circ(t)$, so the exit option is never used. The condition $Q\le x^\circ(0)$ alone is not sufficient. With constant costs, $P^*=\max\{Q,x^\circ\}$.
\end{enumerate}
\end{proposition}

When exit does not bind, Europe's transfer averages the rupture value $M$ of \cref{prop:master}, with weight $\psi$, and the hold-up term $N$, with weight $1-\psi$: Washington's average pain weighted by Europe's agenda power, minus Europe's own pain weighted by Washington's. The hold-up term needs no rupture risk. At $\lambda=0$, \cref{prop:master} gives $\pos{Q}$ whatever $V$ is, while here $P^*=\max\{Q,\beta\hat C_U-(1-\beta)\hat C_E\}$ with constant costs. Hold-up is the fifth channel of \cref{cor:master}(e).

The proposition prices both the sensitivity and the vulnerability of \citet{keohane1977power}, the cost of a cutoff before anyone adjusts and the cost that remains after adjustment. Let baseline disagreement flows be $c_E^0$ and $c_U^0$, which may include Washington's own cost of cutting a loyal Europe's AI access from date zero, and let $x^\circ_0$ be $x^\circ(0)$ under these baseline costs.

\begin{corollary}[Interruption]\label{cor:interrupt}
Let Europe choose, at the start of each round and before the proposer is drawn, whether servicing of US-aligned fabs is interrupted. Interruption costs Europe a flow $\ell$, and costs the United States a flow $(1-h)c_S$ once its buffer is exhausted, where $h$ is the share of that servicing Washington can compel. The buffer is a stock of $t_U$ units of interruption time, depleted only in rounds with interruption. Let $D=\beta(1-h)c_Se^{-(r+\lambda)t_U}-(1-\beta)\ell$. In the frequent-offer limit:
\begin{enumerate}[label=(\roman*)]
\item If exit never binds without interruption (\cref{prop:pain}(iii) with the baseline costs), allowing interruption raises $P^*$ by $(1-\psi)\pos{D}$. Interruption raises $P^*$ if and only if $\beta(1-h)c_Se^{-(r+\lambda)t_U}>(1-\beta)\ell$, that is, if and only if $t_U<\bar t\equiv\ln[\beta(1-h)c_S/((1-\beta)\ell)]/(r+\lambda)$ (with $\bar t=\infty$ if $\ell=0<\beta(1-h)c_S$).
\item With constant baseline costs and the exit option, $P^*=\max\{Q,\,x^\circ_0,\,x^\circ_0+(1-\psi)D\}$. Interruption raises $P^*$ if and only if $(1-\psi)D>\pos{Q-x^\circ_0}$. If exit binds without interruption, the condition of part (i) is necessary but not sufficient.
\item If exit never binds without interruption, as in part (i), the decision is time consistent. At stock $b$, Europe interrupts if and only if $\beta(1-h)c_Se^{-(r+\lambda)b}>(1-\beta)\ell$. Along an interruption path $b$ falls, so the condition keeps holding. If it fails at $t=0$, the stock never moves and it fails in every round.
\end{enumerate}
\end{corollary}

In $D$, vulnerability is the comparison of $(1-h)c_S$ with $\ell$, and sensitivity is the discount $e^{-(r+\lambda)t_U}$ that the buffer buys, which rupture risk and impatience make steep. With constant costs and binding exit, the condition at stock $b$ in part (iii) becomes $(1-\psi)D(b)>\pos{Q-x^\circ_0}$, also time consistent. Timing matters (\cref{app:proofs} proves variants (a) and (c) where exit never binds). If the buffer runs out at a calendar date whatever Europe does, Europe never interrupts before that date (variant (a)). If Europe chooses after an offer is rejected and $\ell>0$, the unique Markov equilibrium has no interruption, because pausing for one round saves Europe $\ell$ and changes nothing that determines the outcome (variant (b)). If it chooses after the proposer is drawn but before offers, it interrupts when it proposes, at no cost on the equilibrium path (variant (c)). Interruption therefore helps only under a rule fixed before offers and kept in force for the dispute. We characterize Markov perfect equilibria; when one side chooses its disagreement action every round, complete-information bargaining generally has many subgame-perfect equilibria \citep{fernandez1991striking, haller1990letter, busch1995perfect}.

The Union's instruments have the wrong timing for such a rule. The Anti-Coercion Instrument allows urgent, immediately applicable acts only to suspend or amend response measures (Article 12(5)) \citep{eu2023aci}, so escalation is slow and committee-gated while de-escalation is fast and in the Commission's hands. An interruption decided this way could be lifted after any US offer, the timing of variant (b), as the Union suspended each countermeasure it armed against Washington (\cref{sec:council}). A credible interruption needs a rule that cannot be suspended after offers, which the Union's current instruments do not supply.

Jurisdiction limits hold-up too. ASML's US subsidiaries and their US staff are US persons \citep{ecfr2026usperson, asml2026subsidiaries}, with offices next to US customer fabs \citep{asml2026locations}. No source we found names the employer of US field-service staff, but on this structure Washington could compel most servicing of US-sited fabs, so $h$ is near one there. ASML's US-sited revenue in 2025 was EUR 4.09 billion \citep{asml2026annual}, so Europe's cost $\ell$ of an interruption confined to US soil would be at most about EUR 0.34 billion a month, and Europe could withhold only new optics, modules and software. For a broad interruption covering the EUV fleet in Taiwan and Korea, $h$ depends on Taipei and Seoul as well as on Washington, and Europe's own cost is larger, about EUR 1.9 billion a month of ASML revenue on a 2025 basis \citep{asml2026annual}. The flows that cannot be commandeered abroad include Zeiss optics and Trumpf drive lasers \citep{asml2026annual, varas2021strengthening}, so on our reading Germany would have to take part in any credible interruption.

Other parameters are poorly measured. The buffer $t_U$ runs from about one month to two years or more. The short end takes at face value an anonymously sourced report that without ASML's spare parts EUV tools ``quickly stop working'' \citep{bloomberg2024killswitch}; the long end allows for stockpiles and for parts and repair capacity seized locally. TSMC's 2025 filing implies US\$3.9 to 4.8 billion a month of wafer revenue from EUV nodes and US-headquartered customers \citep{tsmc2026annual}. With a downstream multiplier of 5 to 12 \citep{varas2021strengthening} and a post-buffer output loss we assume lies between 10 and 70 percent, $c_S$ for a broad interruption runs from about US\$2 billion to US\$40 billion a month, with a central value near US\$13 billion. These are revenues at risk, not welfare losses. With $\beta=1/2$ and central costs, interruption pays at $t_U=0$ only if $h$ is below about $0.85$. With an illustrative $r+\lambda$ of $0.55$ a year, $\bar t$ is about 26 months at $h=0.5$ and about 15 months at $h=0.7$. The answer turns more on $h$ than on $t_U$, and $h$ is the parameter with the least evidence behind it. Since $\bar t$ falls as $h$ rises, ASML's plan to have 50 to 70 percent of service parts repaired locally \citep{asml2026annual} erodes Europe's hold-up option.

Servicing is not Europe's only hold-up instrument. Denying base access, as Spain did for the Iran war (\cref{sec:background}), raises $c_U(t)$, and Washington cannot compel access to another state's territory, so on our reading its $h$ is near zero. One member state decides it, outside the Council, at a cost that includes the conditional part of US posture (\cref{sec:conditionality}). By variant (b), a denial improvised after a rejected offer would never be carried out.

Washington's dependence on ASML bites hardest here. In \cref{prop:master}, $V$ is the loss from China gaining the tools, and it reaches Europe's transfer only through agenda power under rupture risk and through $Q$. In \cref{prop:pain} the relevant loss is $c_S$, what Washington's own fabs lose without ASML, and it raises Europe's transfer without any rupture risk. On US soil $h$ near one leaves little of that loss for Europe to impose; for a broad interruption the answer turns on an $h$ the record does not pin down (\cref{app:scenario} applies this to the scenario's claim that Washington could outlast Europe).
